\section{Methodology}
\label{sec:methodology}

I built my detection pipeline with two primary goals in mind: maximizing my ability to visually expose digital tampering, and ensuring my classification model remained lightweight enough for practical, real-world deployment. To hit both of these targets, I designed a careful, multi-step process that combines targeted image processing with an optimized deep learning framework.

\subsection{Data Curation and Preprocessing}

To train and validate my model, I relied on a carefully selected subset of the FaceForensics++ dataset \cite{faceforensics}. I pulled 10,000 static frames, ensuring I had a perfectly balanced split of 5,000 pristine, authentic images and 5,000 manipulated deepfakes. I also made sure the manipulated group included a wide variety of synthesis techniques. This variety forced my model to learn generalized compression artifacts rather than just memorizing the unique quirks of one specific generative model.

Before feeding any image into the neural network, I passed it through my custom Error Level Analysis (ELA) pipeline, which really forms the heart of my approach. Let an input image be denoted as $I$. I defined the re-compression operation $\mathcal{C}_{q}$, which compresses an image using the JPEG algorithm at a specific quality factor, $q$. For my experiments, I fixed $q$ at 90, as my testing showed this value created the sharpest contrast for compression disparities. I generated the intermediate re-compressed image, $I'$, using the following operation:
\begin{equation}
    I' = \mathcal{C}_{90}(I)
\end{equation}

Next, I calculated the raw ELA difference matrix, $D$, by taking the absolute pixel-by-pixel difference between the original image and its newly compressed counterpart:
\begin{equation}
    D = |I - I'|
\end{equation}

Because the raw difference values in $D$ are usually tiny and hard for a neural network to interpret as primary features, I applied a brightness scaling multiplier. I stretched the pixel intensities to cover the full 0 to 255 color range (8-bit), which maximized the visual contrast. This step effectively strips away the natural visual content of the picture, leaving behind a stark heatmap of the underlying compression grid. In these heatmaps, heavily manipulated areas, such as a generated face, stand out as drastically brighter and noisier than the authentic background.

\subsection{Network Architecture Selection}

Once I had my ELA heatmaps generated, I needed a strong classification backbone. I ultimately went with the EfficientNet-B0 architecture \cite{efficientnet}, primarily because I needed something highly parameter-efficient. While models like ResNet-50 or VGG-16 are undeniably powerful, they are simply too heavy for my operational goals. EfficientNet uses a clever compound scaling method to carefully balance network depth, width, and resolution. This allowed me to achieve state-of-the-art feature extraction without the massive computational cost.

I tailored the standard EfficientNet-B0 to fit my specific binary classification task. First, I removed the original 1000-class ImageNet classification head. In its place, I built a custom top-level module: a global average pooling layer followed by a dense layer with a single linear node. Since I was treating deepfake detection as a straightforward ``Real vs. Fake'' problem, this single node output, combined with a Binary Cross-Entropy with Logits loss function, gave me the cleanest mathematical representation of my objective.

\subsection{The Challenge of Transfer Learning on ELA}

Perhaps the biggest hurdle I faced during my methodology design was navigating the severe domain shift between standard photos and my ELA heatmaps. Like many computer vision researchers, I first tried a standard transfer learning approach. I loaded the ImageNet-pretrained EfficientNet-B0 base, froze the convolutional layers to keep their learned filters intact, and only trained my newly added classification head.

This approach completely failed. The pre-trained base was optimized on ImageNet, a dataset composed of natural photos with smooth gradients, clear edges, and recognizable shapes. Consequently, its lower-level filters were completely unequipped to make sense of my data. ELA heatmaps do not look like natural photos; they resemble abstract, high-frequency, grid-like noise. The frozen filters were simply passing garbage activations up to my classification head.

To overcome this, I pivoted to full network fine-tuning paired with a differential learning rate strategy. Instead of locking down the base network, I opened it up completely. However, I knew the base still contained valuable, generalized edge-detection skills that I did not want to destroy. To protect those skills, I applied a very small learning rate to the base:
\begin{equation}
    \alpha_{base} = 5 \times 10^{-5}
\end{equation}
At the same time, I used a standard, much larger learning rate for the head:
\begin{equation}
    \alpha_{head} = 1 \times 10^{-3}
\end{equation}
